% !TEX program = xelatex
\documentclass[11pt]{article}

\usepackage[a4paper,margin=2.4cm]{geometry}
\usepackage{fontspec}
\usepackage[russian]{babel}
\setmainfont{Times New Roman}
\setsansfont{Arial}
\setmonofont{Consolas}

\usepackage{amsmath,amssymb}
\usepackage{booktabs}
\usepackage{array}
\usepackage{enumitem}
\usepackage{xcolor}
\usepackage{hyperref}
\hypersetup{
  hidelinks,
  pdftitle={Взаимная информация по выборке: MINE и CLUB — лонгрид},
  pdfauthor={Помощник},
  pdfsubject={Разбор тестового задания «Вариационные оценки»}
}

\definecolor{accent}{RGB}{12,74,140}
\setlength{\emergencystretch}{2em}
\setlist{nosep,leftmargin=1.6em}

\newcommand{\nat}{\text{нат}}
\newcommand{\bit}{\text{бит}}
\newcommand{\MI}{I(X;Y)}
\newcommand{\KL}{\operatorname{KL}}
\newcommand{\EE}{\mathbb{E}}

% ---------- Простой «бокс» для примеров (без внешних пакетов) ----------
\newenvironment{exbox}{%
  \par\medskip\noindent
  \begin{minipage}[t]{\linewidth}
  \color{accent}\rule[2pt]{\linewidth}{0.6pt}\color{black}
  \vspace{0.4em}\noindent\small
}{%
  \vspace{0.4em}
  \color{accent}\rule[2pt]{\linewidth}{0.6pt}\color{black}
  \end{minipage}\par\medskip
}

\begin{document}

% ================= Титул =================
\begin{center}
  {\LARGE\bfseries Взаимная информация по выборке:}\\[0.3em]
  {\LARGE\bfseries что это, зачем, и как её оценивают нейросетями}\\[0.8em]
  {\large Лонгрид-разбор тестового задания «Вариационные оценки» (MINE и CLUB)}\\[0.6em]
  {\normalsize\color{accent} Для человека, который знает ML, но впервые слышит об оценке взаимной информации}\\[1.2em]
\end{center}

\begin{exbox}
\textbf{Как читать этот лонгрид.} Раздел 1 — зачем это всё нужно (мотивация, 15~мин). Раздел 2 — термины и формулы с численными примерами (30~мин). Раздел 3 — разбор самого задания: что обучаем, какой таргет, как проверяют (30~мин). Раздел 4 — два готовых файла с кодом, которые можно запустить в Google Colab и потрогать числа руками (1~час). Раздел 5 — план на 8--10 ноября и мини-список литературы. Файлы-спутники: \texttt{00\_baseline\_colab.py} (бейзлайн без нейросетей), \texttt{01\_mine\_minimal.py} (минимальный MINE на PyTorch) и \texttt{02\_club\_minimal.py} (минимальный CLUB на PyTorch).
\end{exbox}

\tableofcontents

% ================= 1. ОТРАСЛЬ =================
\section{Отрасль: что мы делаем и зачем}

\subsection{Где мы находимся}

Задание относится к области \textbf{оценки взаимной информации} (mutual information estimation) — направлению на стыке теории информации, математической статистики и глубокого обучения. Взаимная информация (ВИ, mutual information, MI) $\MI$ — это \emph{универсальная мера статистической зависимости} между случайными величинами $X$ и $Y$. В отличие от коэффициента корреляции Пирсона, который видит только \emph{линейную} связь, ВИ ловит любую зависимость, включая нелинейную.

\begin{exbox}
\textbf{Почему корреляции мало.} Пример 1: $X\sim\mathcal{N}(0,1)$, $Y=X^2$. Связь абсолютно детерминированная, но $\operatorname{corr}(X,Y)=0$ — из-за симметрии распределения квадрат «не коррелирует» со своим аргументом, хотя полностью им определяется. Пример 2: точки, лежащие на окружности: $Y=\sin\Theta$, $X=\cos\Theta$ — корреляция нулевая, зависимость полная. ВИ в обоих случаях больше нуля (а для детерминированной связи — бесконечна). Вывод: если мы хотим мерить «насколько связаны» признаки, изображения или представления, корреляция — слишком грубый инструмент, нужна ВИ.
\end{exbox}

Историческая рамка. ВИ определил Шеннон в 1948 году; долгое время её считали в основном в дискретных задачах (отбор признаков, анализ текстов, филогенетика). В глубоком обучении интерес к ВИ взорвался после серии работ: InfoMax-принцип (Linsker, 1988), InfoGAN (Chen et al., 2016), \textbf{MINE} (Belghazi et al., ICML 2018), InfoNCE (van den Oord et al., 2018) и \textbf{CLUB} (Cheng et al., ICML 2020). Именно MINE и CLUB предлагается реализовать в задании.

\subsection{Зачем оценивать ВИ: сценарии с примерами}

Везде одна и та же механика: есть пара случайных величин, и нужно \emph{число} — насколько сильно они связаны. Разница только в том, \emph{что} потом с этим числом делают.

\begin{itemize}
\item \textbf{Обучение представлений (InfoMax).} Хотим, чтобы представление $z=f(x)$ сохраняло максимум информации о входе $x$: максимизируем $\MI(x;z)$ (Deep InfoMax, 2019). Нужна оценка ВИ, чтобы её дифференцировать и гнать градиентом вверх.
\item \textbf{InfoGAN.} Хотим управлять факторами генерации: максимизируем $\MI(c; G(z,c))$ между кодом фактора $c$ и изображением — сеть учится «проявлять» каждый фактор в отдельной оси латентного пространства.
\item \textbf{Контрастивное обучение (SimCLR, CPC).} Лосс InfoNCE — это, по сути, нижняя граница ВИ между представлением и данными: «тяни похожие пары друг к другу, расталкивай разные». Всё это семейство — родственники MINE.
\item \textbf{Отбор признаков и сжатие.} Выбрать из тысячи признаков те, у которых максимальная $\MI(\text{признак}; \text{таргет})$. Information Bottleneck: сжать $x$ в $z$ так, чтобы $\MI(x;z)$ была мала, а $\MI(z;y)$ велика.
\item \textbf{Unlearning, регуляризация, приватность.} Хотим \emph{уменьшить} зависимость между представлением и чувствительным признаком (пол, возраст, источник данных). Здесь нужно \emph{минимизировать} ВИ — и критически важно, какой границей мы пользуемся (см. ниже).
\item \textbf{Нейронаука.} Оценка информации, которую нейроны несут о стимуле (спайковые данные).
\end{itemize}

\subsection{Ключевая идея направления: нижняя или верхняя}

Точное значение $\MI$ на реальных данных почти никогда не вычислить — нужны плотности $p(x,y)$, $p(x)$, $p(y)$, а у нас лишь выборка пар $\{(x_i,y_i)\}_{i=1}^N$ (картинки, тексты, латентные коды). Оценивать плотности в высокой размерности — отдельная тяжёлая задача (проклятие размерности). Поэтому современный подход — \textbf{вариационные границы}: строим величину, которую можно посчитать прямо по выборке, и которая

\begin{itemize}
\item либо \emph{гарантированно меньше} истинной ВИ (\textbf{нижняя оценка}, lower bound);
\item либо \emph{гарантированно больше} (\textbf{верхняя оценка}, upper bound).
\end{itemize}

Почему важно направление границы? Если вы \emph{максимизируете} ВИ (сжатие, отбор признаков, InfoMax) и оптимизируете нижнюю границу, то истинная ВИ при этом тоже растёт — граница «тащит» её вверх. А вот если вы \emph{минимизируете} ВИ (unlearning) и оптимизируете \emph{нижнюю} границу, то она может падать, а истинная ВИ — расти: сеть просто «сжульничает», обрушив границу. Для минимизации нужна \emph{верхняя} граница. Поэтому в задании написано: нижние оценки — для максимизации (MINE), верхние — для минимизации (CLUB).

\begin{exbox}
\textbf{Пример из жизни.} Вы хотите сделать представление, из которого невозможно восстановить пол пользователя. Вы обучаете сеть минимизировать $\hat I_\text{MINE}$ (нижнюю границу). Сеть находит вырожденного критика, у которого $\hat I_\text{MINE}$ около нуля, — вы довольны. Но настоящая зависимость никуда не делась: при атаке по восстановлению пола точность остаётся высокой. Правильный инструмент для этой задачи — верхняя оценка CLUB: пока она падает, падает и истинная ВИ.
\end{exbox}

Классические не-нейросетевые оценки ВИ (гистограммы, kNN-оценка Кржижанова–Штёгбауэра–Грассбергера, KSG) работают только в малой размерности — их используют как референс, с них мы и начнём в бейзлайне (раздел 4).

% ================= 2. ТЕРМИНЫ =================
\section{Термины и метрики}

\subsection{Энтропия: сколько «сюрприза» в величине}

Для дискретной $X$ с вероятностями $p(x)$:

\[
H(X)=-\sum_x p(x)\log p(x)=\EE[-\log p(X)].
\]

Интуиция: $-\log p(x)$ — это «сюрприз» (чем маловероятнее исход, тем больше), а энтропия — средний сюрприз, мера неопределённости. Если логарифм натуральный, единица измерения — \textbf{нат}; если двоичный — \textbf{бит}. Перевод: $1\ \nat = 1/\ln 2 \approx 1.443\ \bit$.

\begin{exbox}
\textbf{Числа.} Честная монета: $H=-(0.5\log_2 0.5 + 0.5\log_2 0.5)=1$ бит. Кривая монета (орёл с вероятностью $0.9$): $H=-(0.9\log_2 0.9+0.1\log_2 0.1)=0.469$ бит — неопределённости меньше, потому что почти всегда орёл. Честный кубик: $H=\log_2 6\approx 2.585$ бит.
\end{exbox}

Для непрерывных величин используют дифференциальную энтропию $h(X)=-\int p(x)\log p(x)\,dx$; она может быть отрицательной (узкое распределение), но \emph{разности} энтропий, которые нам нужны, — обычные положительные величины. Для гауссианы: $h(\mathcal{N}(\mu,\sigma^2))=\tfrac12\ln(2\pi e\,\sigma^2)$.

\subsection{Совместная и условная энтропии}

Совместная энтропия пары: $H(X,Y)=-\sum_{x,y}p(x,y)\log p(x,y)$ — неопределённость обоих вместе. Условная: $H(X\mid Y)=\sum_y p(y)\,H(X\mid Y=y)$ — сколько в среднем неопределённости остаётся в $X$, если $Y$ известен. Цепное правило: $H(X,Y)=H(Y)+H(X\mid Y)$.

\subsection{Дивергенция Кульбака–Лейблера (KL)}

\[
\KL(P\,\|\,Q)=\sum_x P(x)\log\frac{P(x)}{Q(x)}.
\]

Это мера того, насколько распределение $Q$ «плохо заменяет» $P$. Свойства: $\KL\ge 0$ (неравенство Йенсена); $\KL=0$ тогда и только тогда, когда $P=Q$; \textbf{несимметрична} ($\KL(P\|Q)\neq\KL(Q\|P)$), поэтому это не расстояние. Интуиция кодирования: $\KL(P\|Q)$ — лишние наты, которые вы платите, кодируя данные из $P$ кодом, оптимизированным под $Q$.

\begin{exbox}
\textbf{Числа.} $\KL(\mathrm{Bern}(0.9)\,\|\,\mathrm{Bern}(0.5)) = 0.9\ln(0.9/0.5)+0.1\ln(0.1/0.5)=0.368$ нат, а в обратную сторону $\KL(\mathrm{Bern}(0.5)\,\|\,\mathrm{Bern}(0.9))=0.511$ нат. Одна и та же пара распределений — разные числа: асимметрия в действии.
\end{exbox}

\subsection{Взаимная информация}

Три равносильных определения:

\[
\MI = \KL\big(p(x,y)\,\|\,p(x)p(y)\big)
   = \EE_{p(x,y)}\Big[\log\frac{p(x,y)}{p(x)p(y)}\Big]
   = H(X)-H(X\mid Y).
\]

Словами последней формы: «насколько знание $Y$ снижает неопределённость об $X$». Свойства: симметрична ($\MI=\MI(Y;X)$); неотрицательна; равна нулю \emph{тогда и только тогда}, когда $X$ и $Y$ независимы; $H(X,Y)=H(X)+H(Y)-\MI$.

\begin{exbox}
\textbf{Числа: мороженое и погода.} $X$ — погода ($p(\text{солнце})=0.6$, $p(\text{дождь})=0.4$); $Y$ — «съели мороженое» (при солнце да с вероятностью $0.9$, при дожде — $0.25$). Совместная таблица: $p(\text{солнце},\text{да})=0.54$, $p(\text{солнце},\text{нет})=0.06$, $p(\text{дождь},\text{да})=0.10$, $p(\text{дождь},\text{нет})=0.30$; маргиналы: $p(\text{да})=0.64$, $p(\text{нет})=0.36$. Считаем: $H(X)=0.673$ нат, $H(Y)=0.653$ нат, $H(X,Y)=1.093$ нат, значит $\MI=H(X)+H(Y)-H(X,Y)=0.233$ нат $\approx 0.337$ бит. Погода и мороженое действительно «несут друг о друге информацию», но немного.
\end{exbox}

\subsection{Гауссовский случай — наш рабочий эталон}

Если $(X,Y)$ — совместно гауссовские размерности $d$ с $\operatorname{corr}(x_i,y_j)=\delta_{ij}\rho$ (каждая координата связана только со «своей», с коэффициентом $\rho$), то

\begin{equation}\label{eq:gauss_mi}
\MI = -\frac{d}{2}\log(1-\rho^2).
\end{equation}

Это \emph{главная формула задания}: на такой симуляции истинная ВИ известна, и именно с ней сравнивают оценки. Полезно запомнить её поведение: при $\rho=0$ получаем $\MI=0$ (независимость); при $\rho\to1$ ВИ уходит в бесконечность.

\subsection{Инвариантность к биекциям}

ВИ не меняется при взаимно-однозначных (биективных) преобразованиях аргументов: $\MI(f(X);g(Y))=\MI(X;Y)$ для обратимых $f,g$. Доказательство в одну строку через формулу замены переменных: при $x'=f(x)$ плотность преобразуется как $p(x')=p(x)/|\det J_f(x)|$, и в отношении $p(x,y)/(p(x)p(y))$ якобианы сокращаются, а от отношения плотностей ВИ и зависит.

\begin{exbox}
\textbf{Зачем это в задании.} Инвариантность — сильная проверка честности оценщика: если вы растянули $Y$ преобразованием $y\mapsto\sinh(y)$, информация никуда не делась (преобразование обратимо!), и хороший оценщик должен вернуть то же число. Корреляция Пирсона такой тест проваливает — она меняется при нелинейной деформации.
\end{exbox}

\subsection{«Что мы предсказываем» — главное против заблуждений}

Это \textbf{не} классификация и \textbf{не} регрессия в привычном смысле. В классическом ML есть пары «вход → метка» и сеть предсказывает метку. Здесь \textbf{нет метки} «$\MI=7.3$ ната». Сеть не учится выдавать число $\MI$.

Вместо этого сеть обучается представлять вспомогательный объект, из которого ВИ \emph{вычисляется по формуле}:

\begin{itemize}
\item \textbf{MINE}: сеть-критик $T_\theta(x,y)$ учится приближать \emph{логарифм отношения плотностей} $\log\frac{p(x,y)}{p(x)p(y)}$ — функцию «насколько пара похожа на настоящую». Оценка ВИ — это значение обучающего лосса.
\item \textbf{CLUB}: сеть $q_\theta(y\mid x)$ учится приближать условную плотность $p(y\mid x)$ — это \emph{обычная} задача: по $x$ предсказывать $y$. Оценка ВИ — формула, в которую подставляют уже обученную, замороженную сеть.
\end{itemize}

Истинная формула \eqref{eq:gauss_mi} используется \emph{только} для валидации (сравнить оценку с ответом), в лосс она не входит никогда.

\begin{center}
\small
\begin{tabular}{@{}lccc@{}}
\toprule
Задача & Что предсказывает сеть & Таргет обучения & Где $\MI$ \\
\midrule
Классификация & класс & метка $y$ & — \\
Регрессия & число & значение $y$ & — \\
MINE & скор пары $T(x,y)$ & неявный: $\log\frac{p(x,y)}{p(x)p(y)}$ & = лосс \\
CLUB & распределение $q(y\mid x)$ & настоящий $y$ & формула после обучения \\
\bottomrule
\end{tabular}
\end{center}

\subsection{Как мерить качество оценки}

Раз $\MI$ — это число, оценка $\hat I$ — тоже число, и её качество меряют как у любой оценки: \textbf{смещение} (bias) $\EE[\hat I]-\MI$ и \textbf{дисперсия}. Для отчёта полезны: знак и величина смещения по сетке $\rho$, стандартное отклонение по нескольким сидам, RMSE относительно истинной кривой, и, конечно, \emph{направление границы} ($\hat I\le \MI$ для MINE, $\hat I\ge \MI$ для CLUB).

\subsection{Словарь терминов задания}

\begin{center}
\small
\begin{tabular}{@{}p{5.2cm}p{9.3cm}@{}}
\toprule
Термин & Что значит \\
\midrule
ВИ (MI) & взаимная информация $\MI$; мера зависимости; в задании считается в натах \\
Критик $T_\theta$ & сеть со скалярным выходом на паре $(x,y)$, обучаемая в MINE \\
Положительные пары & $(x_i,y_i)$ из совместного распределения — «настоящие» пары \\
Отрицательные пары & $(x_i,y_{\pi(i)})$: $y$ перемешан в батче — пары «как будто независимые» \\
Перемешивание & сдвиг (shift) или случайная перестановка (\texttt{randperm}) индексов $y$ в батче \\
Лог-сам-эксп & численно устойчивый способ считать $\log\EE[e^T]$: $\operatorname{logsumexp}$ вместо $\exp$ \\
EMA & экспоненциальное скользящее среднее; в MINE — по отрицательному члену, чтобы убрать смещение градиента \\
Нат & единица информации при натуральном логарифме ($1\ \nat\approx1.443\ \bit$) \\
Коллапс $\sigma^2\to0$ & вырождение дисперсии гауссианы в vCLUB, «взрывающее» оценку \\
$\mathcal{O}(N^2)$ / $\mathcal{O}(N)$ & сложность полной/выборочной версии CLUB \\
Абляция & эксперимент с отключённым компонентом метода, показывающий его вклад \\
Биекция & обратимое взаимно-однозначное преобразование; ВИ при нём не меняется \\
Unlearning & удаление нежелательной информации из модели путём минимизации ВИ \\
\bottomrule
\end{tabular}
\end{center}

% ================= 3. ЗАДАЧА =================
\section{Задача: что нужно сделать}

\subsection{Что дано и что требуется}

Дано: выборка пар $\{(x_i,y_i)\}_{i=1}^N$ из неизвестного совместного распределения $p(x,y)$. Плотности недоступны. Требуется: оценить $\MI$ одним из двух нейросетевых оценщиков — \textbf{MINE} (нижняя оценка) или \textbf{CLUB} (верхняя), провалидировать на гауссовской симуляции с известным ответом и написать отчёт с абляциями. Это и есть всё задание.

\subsection{MINE: нижняя оценка}

\paragraph{Теоретическая основа.} Двойственное представление Донскера–Варадхана:

\[
\KL(\mathbb{P}\,\|\,\mathbb{Q}) = \sup_T \Big\{ \EE_{\mathbb{P}}[T] - \log \EE_{\mathbb{Q}}\big[e^{T}\big] \Big\}.
\]

Почему это правда и почему оптимум — именно отношение плотностей? Для любой $T$ определим «перевзвешенное» распределение $R$ с плотностью $dR/dQ = e^{T}/\EE_Q[e^T]$ (это просто нормировка). Тогда прямая подстановка даёт

\[
\EE_P[T]-\log\EE_Q[e^T] = \KL(P\|Q) - \KL(P\|R) \;\le\; \KL(P\|Q),
\]

и равенство достигается ровно когда $P=R$, то есть $e^{T}\propto dP/dQ$. Значит оптимальный критик — это

\[
T^*(x,y)=\log\frac{p(x,y)}{p(x)p(y)} \quad(\text{плюс константа}),
\]

а именно логарифм отношения совместной плотности к произведению маргинальных (точечная взаимная информация, PMI). Суть MINE: \emph{сеть учится логарифму отношения плотностей, не зная самих плотностей}.

\paragraph{Как это считается на батче.} При $\mathbb{P}=p(x,y)$, $\mathbb{Q}=p(x)p(y)$:

\[
\hat I_{\text{MINE}}(\theta)
= \frac1N\sum_{i} T_\theta(x_i, y_i)
- \log\Big(\frac1N\sum_{i} e^{T_\theta(x_i,\, y_{\pi(i)})}\Big)
\;\le\; \MI.
\]

Первый член — средний скор \emph{положительных} пар из совместного распределения. Второй — лог среднего экспоненты скоров \emph{отрицательных} пар, полученных перемешиванием $y$ в батче (сдвиг или \texttt{randperm}): если $x$ и $y$ перемешаны, это ровно выборка из $p(x)p(y)$.

\paragraph{Главный факт: лосс и есть оценка.} Обучение MINE — это максимизация $\hat I_{\text{MINE}}$ градиентным \emph{подъёмом}. Нет отдельной фазы предсказания: после обучения $\hat I_{\text{MINE}}$ — просто значение этого же функционала (на отложенных данных). Поэтому кривая обучения MINE одновременно является кривой оценки.

\paragraph{Что означает каждая рекомендация из задания.}

\begin{enumerate}[label=\arabic*.]
\item \emph{Критик с одним скалярным выходом.} Сеть принимает конкатенацию $[x,y]$ и выдаёт одно число — скор пары.
\item \emph{Отрицательные пары перемешиванием $y$.} Независимые пары негде взять, их имитируют: перемешивание разрушает попарную связь, и $x_i$ с чужим $y_{\pi(i)}$ распределены как $p(x)p(y)$.
\item \emph{Лог-сам-эксп вместо $\exp$.} $e^{T}$ с растущим скором переполняет float. Численно устойчиво: $\log(\frac1N\sum e^{t_i})=\operatorname{logsumexp}(t_1,\dots,t_N)-\log N$. В коде: \texttt{torch.logsumexp(neg) - math.log(N)}.
\item \emph{EMA по отрицательному члену.} Градиент $\log\EE_Q[e^T]$ оценивается по мини-батчу с заметным смещением из-за экспоненциального перевзвешивания (пары с большим скором давят сильнее). Решение из статьи MINE: вести экспоненциальное скользящее среднее величины $\EE_Q[e^T]$ и использовать его в знаменателе лосса.
\end{enumerate}

\paragraph{Ожидаемое поведение.} По построению $\hat I_{\text{MINE}}\le\MI$: оценка должна лежать \emph{ниже} истинной кривой и приближаться к ней по мере обучения. Важно: нижняя оценка может быть \emph{отрицательной} — это нормально (например, при $\rho=0$ критик болтается около нуля, и оценка — около нуля с шумом в обе стороны). На конечной выборке возможен шум, в том числе небольшое превышение истины; систематическое превышение — признак переобучения критика.

\subsection{CLUB: верхняя оценка}

\paragraph{Определение.} При \emph{известной} условной плотности $p(y\mid x)$:

\[
I_{\text{CLUB}} = \EE_{p(x,y)}\big[\log p(y\mid x)\big] - \EE_{p(x)}\EE_{p(y)}\big[\log p(y\mid x)\big].
\]

Почему это верхняя оценка? Добавим и вычтем $\EE\log p(y)$ и соберём KL:

\[
I_{\text{CLUB}} - \MI
= \EE_{p(x)p(y)}\Big[\log\frac{p(y)}{p(y\mid x)}\Big]
= \EE_{p(x)}\Big[\KL\big(p(y)\,\big\|\,p(y\mid x)\big)\Big] \;\ge\; 0.
\]

Итак, $I_{\text{CLUB}}=\MI+\EE_{p(x)}\KL(p(y)\,\|\,p(y\mid x))$: истинная ВИ плюс неотрицательный зазор. Равенство — только при независимости (тогда оба нули).

\paragraph{vCLUB: когда условная плотность неизвестна.} $p(y\mid x)$ заменяют вариационным приближением — например, гауссианой $q_\theta(y\mid x)=\mathcal{N}(y;\ \mu_\theta(x),\ \sigma^2_\theta(x)\cdot I)$, где среднее и дисперсию выдаёт сеть. Подставляя $q_\theta$ в ту же формулу, получают vCLUB. Сеть учат \emph{отдельно}, обычным максимумом правдоподобия: $\frac1N\sum_i\log q_\theta(y_i\mid x_i)\to\max$ (для гауссианы это фактически MSE с обучаемой дисперсией). Потом сеть замораживают и считают оценку:

\[
\hat I_{\text{vCLUB}} = \frac1N\sum_i\log q_\theta(y_i\mid x_i) - \frac1{N^2}\sum_{i,j}\log q_\theta(y_j\mid x_i).
\]

Первый член — насколько правдоподобен \emph{настоящий} $y$ при данном $x$; второй — насколько правдоподобен \emph{случайный} $y$; разница и есть «информативность пары». Важно: vCLUB остаётся верхней оценкой, только если $q_\theta$ хорошо приближает $p(y\mid x)$; при плохом приближении гарантия пропадает.

\begin{exbox}
\textbf{Число для самопроверки (важно!).} В гауссовской постановке задания $p(y\mid x)=\mathcal{N}(\rho x,\ (1-\rho^2)I)$ — то есть условная плотность \emph{сама} гауссова, и сеть может выучить её точно. Если $q_\theta=p(y\mid x)$, то $I_{\text{CLUB}} = d\rho^2/(1-\rho^2)$. Это \emph{больше} истинной ВИ из \eqref{eq:gauss_mi}: например, $d=1$, $\rho=0.8$ даёт $I_{\text{CLUB}}=1.78$ нат против $\MI=0.51$ нат. Вывод: vCLUB в этой постановке — \emph{не точная} оценка, а честная верхняя граница с ненулевым зазором, и её значение $d\rho^2/(1-\rho^2)$ можно предсказать заранее. Если ваша vCLUB-оценка близка к этому числу — сеть выучилась хорошо; если она ещё выше — что-то не так (обычно коллапс $\sigma^2$).
\end{exbox}

\paragraph{Что означает каждая рекомендация из задания.}

\begin{enumerate}[label=\arabic*.]
\item \emph{Учить только на лог-правдоподобии.} Если дифференцировать полную оценку, градиент потечёт через двойную сумму $\mathcal{O}(N^2)$ — медленно и шумно, а главное — сеть начнёт «оптимизировать под оценку», ломая приближение. Правильно: фаза 1 — обычная регрессия $y\mid x$ (MLE), фаза 2 — заморозить и посчитать оценку без градиентов.
\item \emph{Выборочная версия $\mathcal{O}(N)$.} Второй член можно оценивать не по всем $N^2$ парам, а по $N$ парам $(x_i, y_{j(i)})$ со случайным $j(i)$: получится несмещённая оценка того же выражения ценой чуть большей дисперсии.
\item \emph{Коллапс $\sigma^2\to0$.} При переобучении сеть уменьшает дисперсию: $\log q_\theta(y\mid x)=-\tfrac{d}{2}\log(2\pi\sigma^2)-\frac{\|y-\mu\|^2}{2\sigma^2}$, и при $\sigma^2\to0$ первый член уходит в $+\infty$ — оценка «взрывается» вверх. Лечение: параметризовать $\log\sigma^2$ вместо $\sigma^2$, клиппинг $\sigma^2\ge\sigma^2_{\min}$, общая (не зависящая от $x$) дисперсия, early stopping.
\end{enumerate}

\subsection{Таргеты: сводка «кто что учит»}

\begin{center}
\small
\begin{tabular}{@{}p{3.1cm}p{5.6cm}p{5.8cm}@{}}
\toprule
 & MINE & CLUB (vCLUB) \\
\midrule
Сеть & критик $T_\theta(x,y)\to\mathbb{R}$ & $q_\theta(y\mid x)$: гауссиана с $\mu_\theta,\sigma^2_\theta$ \\
Таргет обучения & нет явного; неявный — $\log\frac{p(x,y)}{p(x)p(y)}$ & настоящий $y$ \\
Лосс & $\hat I_{\text{MINE}}$ (максимизация) & $-\frac1N\sum\log q_\theta(y_i\mid x_i)$ (NLL) \\
Когда читают оценку & значение лосса (обучение $=$ оценивание) & после обучения, формулой по замороженной сети \\
Направление & $\hat I\le\MI$ (нижняя) & $\hat I\ge\MI$ при хорошем $q$ (верхняя) \\
Истинная $\MI$ в лоссе & нет & нет \\
\bottomrule
\end{tabular}
\end{center}

Повторим главное: \textbf{ни в одном методе сеть не учится предсказывать число $\MI$}. Истинная ВИ из формулы \eqref{eq:gauss_mi} — внешний «учитель», который появляется только при валидации, чтобы сверить ответ.

\subsection{Валидация по заданию}

\paragraph{Генератор.} $(X,Y)$ — многомерное гауссово размерности $d$ с $\operatorname{corr}(x_i,y_j)=\delta_{ij}\rho$: генерируем $x\sim\mathcal{N}(0,I)$ и $y=\rho x+\sqrt{1-\rho^2}\,\varepsilon$, где $\varepsilon\sim\mathcal{N}(0,I)$ независимо. Истинное значение — формула \eqref{eq:gauss_mi}. Меняя $\rho$, получаем сетку значений $\MI$ в диапазоне 2--10 нат.

\paragraph{Как выбрать $d$ и $\rho$.} При $d=1$ даже $\rho=0.99$ даёт всего $\approx2.6$ нат, поэтому удобно брать $d=10$ и сетку, например:

\begin{center}
\small
\begin{tabular}{@{}lcccccc@{}}
\toprule
$\rho$ & 0.45 & 0.55 & 0.65 & 0.75 & 0.82 & 0.92 \\
\midrule
$\MI$, наты ($d=10$) & 1.13 & 1.80 & 2.75 & 4.13 & 5.58 & 9.37 \\
\bottomrule
\end{tabular}
\end{center}

\paragraph{Три обязательные проверки.}

\begin{enumerate}[label=\arabic*.]
\item \emph{Направление границы}: MINE должен давать $\hat I\le\MI$, CLUB — $\hat I\ge\MI$, по всей сетке $\rho$.
\item \emph{Независимость}: при $\rho=0$ оценка должна быть $\approx0$. Толкование «$\approx$»: в пределах шума (MINE может быть отрицательным, у классических оценщиков вроде KSG остаётся небольшой конечновыборочный шум) — главное, что оценка на порядки меньше, чем при $\rho>0$.
\item \emph{Инвариантность}: применить биекции (например, $y\mapsto\sinh(y)$ покомпонентно или обратимое линейное $x\mapsto Ax$) — оценка должна остаться той же. Ловушка для CLUB: после нелинейной биекции данные перестают быть гауссовыми, и гауссова $q_\theta$ приближает $p(y'\mid x)$ хуже — оценка деградирует, это стоит обсудить в отчёте как ограничение.
\end{enumerate}

\paragraph{Ловушки, которые проверяющий увидит (или не увидит) сразу.}

\begin{itemize}
\item MINE-оценка может быть отрицательной — это нормально, в отчёте так и напишите.
\item Переобученный критик систематически завышает MINE — держите оценку на отложенной выборке.
\item Для MINE важен размер батча: отрицательные пары берутся из батча, и качество границы зависит от него — это готовая абляция.
\item Несколько сидов и полосы $\pm$std на всех графиках — иначе «результат» не считается воспроизводимым.
\item Гиперпараметры (скорость, ширина сети, число шагов) — явно в отчёте.
\end{itemize}

\subsection{Что сдавать и как оценивают}

Критерий отбора из задания: \emph{правильная реализация, подробный отчёт, абляционные эксперименты}. Минимальный состав сдачи: код (генератор + оценщик + скрипты валидации), график $\hat I$ против истинной кривой с полосами по сидам, три проверки выше, 2--3 абляции (например: log-sum-exp вкл/выкл, EMA вкл/выкл, сдвиг vs \texttt{randperm}, архитектура критика — линейная vs MLP; для CLUB: полная vs выборочная версия, клиппинг $\sigma^2$, гауссова $q$ на не-гауссовых данных), и текст отчёта. \emph{Совет: один метод, сделанный честно и до конца, сильнее двух недоделанных.} CLUB надёжнее (обычная регрессия + формула, и вы заранее знаете, к какому числу должна сойтись оценка); MINE эффектнее (нижняя граница с EMA), но капризнее.

% ================= 4. БЕЙЗЛАЙН =================
\section{Бейзлайн: потрогать числа руками}

\subsection{Зачем бейзлайн}

Прежде чем писать нейросетевой оценщик, полезно увидеть саму величину: сгенерировать данные, посчитать истинную ВИ по формуле \eqref{eq:gauss_mi}, оценить её классическими методами и посмотреть на числа и графики. Файл \texttt{00\_baseline\_colab.py} делает ровно это без единой нейросети (numpy + scipy + matplotlib, всё есть в Colab). Запуск в Colab:

\begin{verbatim}
exec(open('00_baseline_colab.py', encoding='utf-8').read())
\end{verbatim}

\subsection{Что внутри файла}

\begin{enumerate}[label=\arabic*.]
\item \textbf{Генератор} \texttt{sample\_gauss(d, rho, n)} — ровно модель из задания: $y=\rho x+\sqrt{1-\rho^2}\,\varepsilon$.
\item \textbf{Истинная ВИ} \texttt{true\_mi(d, rho)} — формула \eqref{eq:gauss_mi}.
\item \textbf{Оценка №1: гистограммная} (\texttt{mi\_hist}, только $d=1$). Простейшая идея: нарезать $x$ и $y$ на бины, оценить частоты $p(x,y)$, $p(x)$, $p(y)$ и подставить в определение ВИ. Преподаёт два урока: (а) оценка смещена — при малом $N$ и многих бинах завышает, при мало бинов — занижает; (б) клеток $B^{2d}$ — проклятие размерности, поэтому только $d=1$.
\item \textbf{Оценка №2: KSG} (\texttt{mi\_ksg}, kNN-оценка Кржижанова–Штёгбауэра–Грассбергера). Классический непараметрический оценщик: для каждой точки ищут расстояние до $k$-го соседа в совместном пространстве, считают соседей в маргинальных, и комбинируют через дигамма-функцию. Необязателен для задания, но это отличная «эталонная линия» для отчёта.
\item \textbf{График 1}: KSG против истинной кривой по сетке $\rho$ — должны почти совпасть (с полосами по сидам).
\item \textbf{График 2} (тест независимости): $\rho=0$, гистограмма оценок по 30 сидам — должны скучиться около нуля.
\item \textbf{График 3} (инвариантность): $y\mapsto\sinh(y)$ покомпонентно — KSG должен вернуть то же значение, а корреляция Пирсона изменится. Числа печатаются в консоль.
\item \textbf{График 4} (нелинейная зависимость): $y=x+\varepsilon$ против $y=x^2+\varepsilon$ — корреляция во втором случае около нуля, а ВИ заметно больше нуля. Наглядная демонстрация, зачем вообще ВИ.
\end{enumerate}

\subsection{Чеклист понимания после прогона}

Если вы можете ответить на пять вопросов — можно переходить к нейросетям. (1) Почему кривая $\MI(\rho)$ растёт и уходит в бесконечность при $\rho\to1$? (2) Почему гистограммная оценка завышает при большом числе бинов? (3) Почему KSG при $\rho=0$ даёт не ровно ноль, а значение с небольшим шумом вокруг нуля? (4) Почему после $\sinh$ корреляция изменилась, а ВИ — нет? (5) Почему ВИ видит зависимость $y=x^2$, а корреляция — нет?

\subsection{Следующий шаг: минимальный MINE}

Файл \texttt{01\_mine\_minimal.py} — минимальный рабочий MINE на PyTorch (в Colab уже установлен): критик-MLP на конкатенации $[x,y]$, лосс с \texttt{logsumexp} и $\log N$, перемешивание через \texttt{randperm}, EMA-опция, обучение 400 шагов, оценка на отложенной выборке, график $\hat I$ против истинной кривой и кривая обучения. В файле есть два переключателя-абляции: \texttt{use\_ema} (вкл/выкл EMA) и замена MLP на линейный критик $w\cdot(x^\top y)$ — запустите оба варианта и посмотрите, как линейный критик систематически занижает оценку (он не может представить оптимальный $T^*$, который в гауссовском случае является квадратичной формой). Это готовый прототип для настоящей реализации: добавьте сиды, сетку $\rho$ под 2--10 нат ($d=10$), проверку инвариантности и абляции.

\subsection{Ещё один шаг: минимальный CLUB}

Файл \texttt{02\_club\_minimal.py} — минимальный рабочий CLUB: сеть $q_\theta(y\mid x)$ в виде диагональной гауссианы (сеть выдаёт $\mu$ и $\log\sigma^2$, отсюда идёт клиппинг от коллапса $\sigma^2\to0$), фаза~1 — обычный MLE по совместным парам, фаза~2 — оценка по замороженной сети с переключателем «полная $\mathcal{O}(N^2)$ / выборочная $\mathcal{O}(N)$». На графике рядом с истинной кривой нарисована \emph{предсказанная} граница $d\rho^2/(1-\rho^2)$ — если обученная vCLUB близка к ней, значит $q$ выучена хорошо; если ниже истинной ВИ — $q$ недоучена, и гарантия верхней оценки теряется. Рекомендация: если времени мало, в качестве основного метода сдавайте CLUB (проверяемо и устойчиво), а MINE оставьте для сравнения в отчёте.

% ================= 5. ПЛАН =================
\section{План до 10 ноября}

Ниже — реалистичный разрез по дням (итого 12--15 часов). Принцип приоритетов: сначала \emph{корректность метода и направление границы}, потом три проверки, потом абляции, в самом конце — полировка отчёта. Если время жмёт, режьте снизу списка (сначала число абляций), но никогда не жертвуйте корректностью метода.

\begin{center}
\small
\begin{tabular}{@{}p{2.6cm}p{9.6cm}p{2.3cm}@{}}
\toprule
День & Что делать & Бюджет \\
\midrule
Вс 8.11 & \textbf{Понять и пощупать.} Прочитать лонгрид (§1--2) + «Visual Information Theory» К.~Олаха для интуиции; перечитать §3 и исходный PDF задания; запустить \texttt{00\_baseline\_colab.py} в Colab, покрутить $d$, $\rho$, $N$, бины, $k$ и закрыть чеклист из §4.3; выбрать метод (рекомендация: CLUB — надёжнее) и запустить \texttt{02\_club\_minimal.py}, сверив оценку с предсказанной линией $d\rho^2/(1-\rho^2)$ & 2--3 ч \\
\hline
Пн 9.11 & \textbf{Сделать по-настоящему.} Довести оценщик до требований задания: log-sum-exp и EMA (MINE) или $\log\sigma^2$-параметризация + клиппинг и выборочная $\mathcal{O}(N)$-версия (CLUB); сетка $\rho$ под 2--10 нат ($d=10$); 3--5 сидов; оценка на отложенных данных. Провести три проверки валидации с графиками. Сделать 2--3 абляции из списка §3.6 & 5--6 ч \\
\hline
Вт 10.11 & \textbf{Отчёт и сдача.} Написать отчёт (постановка → метод → протокол → результаты → абляции → выводы), вставить графики с полосами по сидам, таблицу гиперпараметров. Финальный прогон всего с нуля (проверка воспроизводимости), контроль направлений границ и $\rho=0$, упаковка и отправка & 3--4 ч \\
\bottomrule
\end{tabular}
\end{center}

\subsection{С чего начать изучение темы}

\begin{enumerate}[label=\arabic*.]
\item Этот лонгрид, разделы 1--2 — термины и интуиция (45 мин).
\item Chris Olah, «Visual Information Theory» (colah.github.io) — картинки про энтропию, KL и ВИ (20 мин).
\item Раздел 3 + исходный PDF задания — постановка и ловушки (30 мин).
\item Бейзлайн \texttt{00\_baseline\_colab.py} — числа руками (1 ч).
\item Статья выбранного метода, только введение и формулы, без доказательств: MINE — arXiv:1801.04062 (§1--3); CLUB — arXiv:2006.12013 (§1--3) (30--40 мин).
\item Реализация (файлы-стартеры) и валидация — остальное время.
\end{enumerate}

\begin{exbox}
\textbf{Резюме для тех, кто читает по диагонали.} Задание — оценить взаимную информацию пары случайных величин по выборке без плотностей. Сеть не предсказывает число ВИ: MINE учит критика — логарифм отношения плотностей, и его лосс \emph{является} нижней оценкой; CLUB учит обычную регрессию $q(y\mid x)$, а верхнюю оценку считают формулой по замороженной сети. Валидация — гауссовская симуляция с известным ответом $-\frac d2\log(1-\rho^2)$: проверить направление границы, нулевую оценку при независимости и инвариантность к биекциям. Сдать: код, отчёт, абляции. До 10-го реально успеть, если начать сегодня с раздела 4.
\end{exbox}

\end{document}
